% explainer-source-hash: f4056f58dd7b4a388ee50f223c77b4315cabd8562bb0f0a0f393b8a5e9d6e423
% explainer-source-commit: 2b8bb12994f7dc6f75fe0d99b2cff9e1de273266
% explainer-generated: 2026-07-16
% Regenerate with /explain-all (see WIRING.md). Do not edit this stamp by hand:
% it is how the toolchain knows whether this document still matches the code.
\documentclass[11pt,a4paper]{article}
\input{preamble}

\begin{document}

\begin{center}
{\LARGE\bfseries\color{inkblue} Agency Blog SaaS: Demo\par}
\vspace{4pt}
{\large\color{accent} Brief\par}
\vspace{10pt}
{\small A short, accurate account of what this project is, what it proves, and where it
falls short.\par}
\end{center}
\vspace{6pt}

\begin{honest}{The framing, first, because everything depends on it}
This repository is \textbf{not} the product. It is a static web page that draws pictures of
a product.

The real system, \file{agency-blog-saas}, is a production Django SaaS behind
digitalise.agency with real clients, real billing and real AI generation. Its source is
private and stays private: it is the agency's live product, not a codebase that could be
scrubbed and released. This demo exists so the product's screens can be clicked by someone
who will never be shown the real repository.

Nearly everything on the page is \textbf{fabricated}: the credit balance, the client, the
updates, the generation itself. Nothing is computed, nothing is sent anywhere. The one
exception is the blog library, which reproduces a real client's already-public articles, and
the page labels that distinction in four separate ways.
\end{honest}

\section{What it is}

One HTML page, one CSS file, and two JavaScript files totalling 216 lines of behaviour plus
150\,KB of stored data. No framework, no build step, no backend, no dependencies, no
secrets, and no network calls at runtime. Deployment is two lines of
\file{netlify.toml} saying ``publish the folder as it is''. To run it, open
\file{index.html}.

Four tabs mirror the real product's screens, and a fifth hidden panel reads an article in
place:

\begin{itemize}
\item \textbf{AI Blog Generator}: a form, an Autofill button, and a Generate button that
      runs a fixed 2.9-second sequence of status messages, subtracts 100 from a number, and
      produces a post that was already sitting in the page.
\item \textbf{Blog Library}: fifteen real, live articles from a real client's public blog.
\item \textbf{Subscriptions}: three real, published price tiers and a redemption box that
      always answers that it is a demo.
\item \textbf{Client Updates}: three lines for an openly invented company.
\item \textbf{Reader}: opens any card into a full article with hero image, prose, tables,
      citations, references, and a link back to the original.
\end{itemize}

\section{The idea that makes it worth reading}

The interesting problem here is not technical. A fabricated demo is a fiction, and the whole
design question is how to keep a visitor from mistaking the fiction for the fact.

The project's answer is that data on the page comes in three tiers, and each is labelled by
a different mechanism.

\begin{figure}[H]
\centering
\resizebox{\textwidth}{!}{
\begin{tikzpicture}
\node[box, fill=warnred!10, draw=warnred, text width=4.0cm, minimum height=2.5cm,
      anchor=north west] (t1) at (0,0)
  {\textbf{Fabricated}\\[3pt]\scriptsize 1{,}200 credits, ``Aurora\\ Fitness Co.'', its
   updates,\\ the generation timeline};
\node[box, fill=softgrey, text width=4.0cm, minimum height=2.5cm,
      anchor=north west] (t2) at (4.6,0)
  {\textbf{Real, public,}\\\textbf{not client-specific}\\[3pt]\scriptsize the price tiers,
   the\\ stack line, the test line,\\ davonex.com's design tokens};
\node[box, fill=okgreen!10, draw=okgreen, text width=4.0cm, minimum height=2.5cm,
      anchor=north west] (t3) at (9.2,0)
  {\textbf{Real client work,}\\\textbf{already public}\\[3pt]\scriptsize 15 davonex.com
   articles:\\ real text, images, dates,\\ 177 references};

\node[extbox, text width=4.0cm, anchor=north west] (x1) at (0,-2.9)
  {\scriptsize A permanent red banner,\\ plus the client's name is\\ literally stored as
   ``(invented,\\ illustrative client)''.};
\node[extbox, text width=4.0cm, anchor=north west] (x2) at (4.6,-2.9)
  {\scriptsize Safe because it was public\\ before the demo existed.\\ Reused verbatim, not\\
   approximated.};
\node[extbox, text width=4.0cm, anchor=north west] (x3) at (9.2,-2.9)
  {\scriptsize A provenance note and a\\ link to the original URL at\\ the foot of every
   article.};
\draw[flow] (t1.south) -- (x1.north);
\draw[flow] (t2.south) -- (x2.north);
\draw[flow] (t3.south) -- (x3.north);

\node[box, draw=accent, text width=13.2cm, anchor=north west] at (0,-5.2)
  {\textbf{The unifying rule:} anything fabricated is a \emph{duplicate} of something real,
   stamped \code{Demo} with a live countdown. The demo never invents a post; it copies one
   and labels the copy.};
\end{tikzpicture}}
\caption{Three provenances on one page, and the machinery that keeps them apart.}
\end{figure}

Shipping a real client's genuine, shipped work beside clearly-labelled fakes is an unusual
choice and a defensible one: the visitor sees real output and a working interaction in one
place, with a checkable link on every article.

\section{Architecture}

\begin{figure}[H]
\centering
\resizebox{\textwidth}{!}{
\begin{tikzpicture}
\node[box, text width=2.6cm, minimum height=1.5cm] (html) at (0,0)
  {\file{index.html}\\\scriptsize 5 empty\\ panels};
\node[box, text width=3.5cm, minimum height=1.5cm] (data) at (6.0,0)
  {\file{demo-data.js}\\[2pt]\scriptsize \code{REAL\_DAVONEX\_POSTS}\\
   \code{GENERATOR\_PRESETS}\\ \code{DEMO}};
\node[box, text width=2.4cm, minimum height=1.5cm] (app) at (11.8,0)
  {\file{app.js}\\\scriptsize behaviour};
\node[store, text width=1.5cm] (dom) at (17.9,0) {live\\ DOM};
\draw[flow] (html) -- node[lbl, above=1pt] {script 1} (data);
\draw[flow] (data) -- node[lbl, above=1pt] {script 2\\ (order is\\ load-bearing)} (app);
\draw[flow] (app) -- node[lbl, above=1pt] {the 4 renderers,\\ on\\
  \code{DOMContentLoaded}} (dom);
\node[box, draw=warnred, fill=warnred!8, text width=5.6cm] (net) at (6.0,-3.2)
  {\textbf{There is no network layer.}\\\scriptsize No \code{fetch}, no external URL, no font
   CDN.\\ Verified: \textbf{0} non-local requests.};
\draw[dflow] (net.north) -- (data.south);
\end{tikzpicture}}
\caption{The complete architecture. Two script tags, three globals, four render functions.}
\end{figure}

Each file has exactly one job: HTML is structure, \file{demo-data.js} is data,
\file{app.js} is behaviour, \file{styles.css} is appearance. Because both scripts are plain
(non-module) scripts, their top-level constants share one global scope, and \file{app.js}
reads \file{demo-data.js}'s constants by name. That makes the order of two lines in the HTML
a real dependency that neither file records.

The one design decision worth defending is that \code{GENERATOR\_PRESETS} is \emph{derived}
from the posts rather than authored:

\begin{lstlisting}
const GENERATOR_PRESETS = REAL_DAVONEX_POSTS.map(p => ({
  slug: p.slug, topic: p.title, description: p.excerpt,
  keywords: p.title.toLowerCase().replace(/[():]/g, "")
              .split(" ").filter(w => w.length > 3).slice(0, 6).join(", ")
}));
\end{lstlisting}

\begin{keyidea}{Why that matters}
If Autofill wrote one thing into the form and Generate produced something unrelated, the
demo would be quietly lying about the one property a generator must have: that the output
follows from the input. Deriving each preset \emph{from} the post it will produce makes that
link structural rather than clerical. Preset $i$ carries post $i$'s slug; Autofill records
it; Generate looks it up. They cannot drift apart because there is nothing to keep in sync.

(The README adds that this mirrors the real pipeline's deterministic-per-brief behaviour.
That is a claim about the private system and cannot be checked from this repository.)
\end{keyidea}

Two other techniques are worth naming. \textbf{Event delegation}: the library's cards are
destroyed and rebuilt every second by a countdown, so no listener is attached to a card; one
listener sits on the container and uses \code{e.target.closest("[data-open-post]")} to find
which card was hit. \textbf{Single-slot state}: \code{libraryDemo} is one object or
\code{null}, never a list, so generating twice replaces rather than accumulates.

\section{Verified, not assumed}

The page was driven in a real headless Chrome and the data evaluated in Node.

\begin{table}[H]
\centering
\small
\begin{tabular}{@{}p{0.52\textwidth}p{0.42\textwidth}@{}}
\toprule
\textbf{Claim} & \textbf{Result} \\
\midrule
No network calls at runtime & \textbf{True}: 0 non-local requests \\
No JavaScript errors across the full tour & \textbf{True}: 0 errors \\
15 real posts, unique slugs and titles, images on disk & \textbf{True} \\
Generate adds one card, credits 1{,}200 to 1{,}100 & \textbf{True} \\
The autofilled post is the generated post & \textbf{True} \\
Article renders with prose, table, 11 references & \textbf{True} \\
\midrule
Citation links resolve & \textbf{False}: 21 links, 21 dead \\
A typed topic is respected & \textbf{False}: returns an unrelated post \\
Credits have a floor & \textbf{False}: renders $-50$ \\
The ticker stops when the card expires & \textbf{False}: runs forever \\
\midrule
132/132 tests, 56\% coverage & \textbf{Not verifiable here} \\
CSS tokens copied 1:1 from davonex.com & \textbf{Not verifiable here} \\
\bottomrule
\end{tabular}
\end{table}

\section{What does not hold up}

\begin{honest}{The banner's central claim is false as written}
The red banner says \emph{``nothing shown is a real client's data''}, and the footer repeats
it. But the Blog Library is fifteen real posts from a real client, and the project's own
README, CREDITS and provenance notes say so proudly.

The intended meaning is obviously ``no \emph{private} client data'', and everything shown is
genuinely public. But as written, the page's most prominent honesty statement contradicts
the page's most interesting honest feature, handing a careful reader a reason to distrust
every other label. ``No private client data; the Blog Library is a real client's
already-public content, labelled as such'' would be both accurate and stronger.
\end{honest}

\begin{honest}{Generating without autofilling returns someone else's post}
\code{const slug = currentPresetSlug || GENERATOR\_PRESETS[0].slug;}

Type your own topic, ignore Autofill, hit Generate, and the fallback fires. Verified: typing
\code{My Own Unrelated Topic} produces \code{Done: Track AI Search Sales: E-Commerce
Attribution Guide (2026) added to your Blog Library}. Nothing warns the visitor that their
input was discarded, and the fields are neither required nor read-only.

The most natural first interaction on the page silently produces an unrelated article and
asserts it as the result: the exact thing the preset-derivation design was built to prevent,
one branch away from all that care. This is the only unlabelled untruth on the page.
\end{honest}

\begin{honest}{Every citation link in every article is dead}
The real article bodies are full of \code{<sup class="citation"><a href="\#ref-1">} markers.
The reader renders references as a plain \code{<ol>} with no \code{id} attributes, and no
\code{id="ref-N"} target exists anywhere in the stored data either. Measured on one opened
post: \textbf{21 citation links, 21 dead}.

This is the price of copying rendered HTML out of a site into a different template. It is a
one-line fix, and worth doing, because the citations are the part of the demo that shows
real, sourced work.
\end{honest}

\begin{honest}{The countdown's interval leaks forever}
\code{setTimeout(() => \{ libraryDemo = null; renderLibrary(); \}, ...)} nulls the object
that holds the interval handle, so \code{clearInterval} can never be called for it. Verified
by counting DOM rebuilds: 4 per 3 seconds while the card lives, and \textbf{4 per 3 seconds
after it expires}, permanently. Each tick rebuilds the entire fifteen-card grid, images and
all, via \code{innerHTML}, to update one \code{mm:ss} string.

The replace path does clear both handles correctly; only the expiry path forgets. One line
fixes it. The deeper point, that a countdown should update a text node rather than
re-render a grid, is the clearest ``I would do it differently'' the codebase offers.
\end{honest}

\begin{honest}{Smaller, but real}
\begin{itemize}
\item \textbf{Credits go negative.} \code{creditsBalance -= 100} has no floor, in a demo
      whose subject is metering. Thirteen clicks reach $-100$. One \code{if} would turn a
      bug into a demonstration of the feature.
\item \textbf{\code{DEMO.facts} is dead code}, never read anywhere, while the same three
      figures are hard-coded again in the footer of \file{index.html}. Two copies, one
      wired to nothing, in a project whose premise is traceable claims.
\item \textbf{The derived keywords drop the subject.} The length filter deletes ``Ads'',
      ``SEO'', ``AEO'', ``CAC'' and ``vs'', so the Google Ads post's keywords contain
      neither ``ads'' nor ``AEO''; stopwords like ``your'' and ``what'' survive; and the
      punctuation strip misses \code{?}, shipping the literal keyword \code{2026?} on
      screen.
\item \textbf{Two advertised capabilities are absent.} The hero names per-client editor
      roles and PDF export; neither appears anywhere in the tour.
\item \textbf{The provenance names a different system.} \file{CREDITS.md} says davonex.com
      is a client of \file{blog-posting-pipeline}, a \emph{separate} private project, while
      the page's prose says only ``the real production pipeline'' inside a demo titled
      Agency Blog SaaS. A visitor will reasonably conclude these posts came from this
      product. Worth settling before an interview.
\item \textbf{The escaping is theatre.} \code{escapeHtml} is applied diligently, but every
      string it touches is a hard-coded literal and no user input ever reaches the DOM. The
      page is safe because it has no untrusted input, not because it escapes. Good habit,
      accurate description needed.
\end{itemize}
\end{honest}

\section{The honest bottom line}

\begin{keyidea}{What this demonstrates, and what it does not}
It demonstrates that its author can build a small, clean, dependency-free static page, hold
a strong self-imposed constraint (no network calls) that survives measurement at zero
non-local requests, and think unusually carefully about how to label a fiction so nobody
mistakes it for fact.

It contains \textbf{no evidence} about the real product: not whether generation works, not
output quality, not prompting, not metering or refunds, not billing, not roles, not PDF
export, not performance, not security, not scale, and not whether the tests it cites exist.

The strongest defensible sentence is: ``This is a puppet, I built the puppet carefully, and
I labelled it as a puppet.'' Everything in this document supports that. Nothing in it
supports a claim about the real system's quality, and the demo is at its best when it does
not try to make one.
\end{keyidea}

\section*{Glossary}
\addcontentsline{toc}{section}{Glossary}

\begin{description}[leftmargin=0pt, style=nextline, itemsep=3pt]
\item[Static site] A site made of files served exactly as stored, with nothing generated per
visitor. The opposite is a dynamic site, where a server program (Django, say) builds each
page from a database. The real system is dynamic; this demo is static, which is why it can
only imitate dynamism by scripting it in the browser.
\item[DOM] The browser's live in-memory tree of the page. HTML is the text you wrote; the
DOM is the object the browser built from it and actually draws. JavaScript changes the page
by changing the DOM.
\item[\code{innerHTML}] The blunt way to change the DOM: assign a string of HTML and the
browser discards everything inside the element and re-parses it. Simple to write; it
destroys and recreates elements every time, and it executes whatever markup the string
contains.
\item[Event delegation] Attaching one listener to a container instead of one per child, then
using \code{closest()} to find which child was hit. Necessary here because the cards are
rebuilt every second and per-card listeners would die with them.
\item[Slug] A lowercase, hyphenated, URL-safe form of a title. Here it is both a post's
identity inside the demo and the real path on the real site, which is how the provenance
link is built by concatenation rather than stored.
\item[\code{setTimeout} / \code{setInterval}] Ask the browser to call a function once after
a delay, or repeatedly at an interval. Both return a handle; losing the handle means the
timer can never be cancelled, which is exactly the bug above.
\item[XSS and escaping] Cross-site scripting: if attacker-controlled text is pasted into
HTML, the browser cannot tell it from your markup and will run any \code{<script>} in it.
Escaping turns \code{<} into \code{\&lt;} so it is drawn rather than parsed.
\item[CSS custom properties] Variables like \code{--blue: \#136BEE} declared on \code{:root}
and read with \code{var(--blue)}, so the palette lives in one block instead of being
scattered as literal hex codes.
\item[Netlify, \code{publish = "."}] A host that serves a git repository as a website. Most
projects need a build command first; this one has no build, so the config says only
``publish the folder as it is''.
\item[PolyForm Strict 1.0.0] A source-available licence: read and evaluate the code, but do
not use, modify or redistribute it. The right choice here, because the repository contains a
real client's copyrighted articles that are not the licensor's to give away freely.
\end{description}

\end{document}
